\documentclass[11pt,a4paper]{article}
\usepackage[margin=25mm]{geometry}
\usepackage{amsmath,amssymb,amsthm,hyperref}
\newtheorem{lemma}{Lemma}
\newcommand{\gen}[1]{\langle#1\rangle}
\title{A regular four-point orbit at the crossing anchor}
\author{Human Compression project: bounded L044 proof component}
\date{}
\begin{document}\maketitle
Retain the accepted unique-crossing string C-group setting in
$G\leq A_n$, $n\geq12$, with all maximal omissions intransitive. The two
$G_i$-orbits $A,B$ have two blocks of sizes $k,l$. Assume
$\alpha=0$, $\beta=1$, $\gamma=0$, and that the crossing
endpoint $a\in A$ lies in the regular four-point orbit $D$
of $M\cong C_2^2$. All other $M$-orbits are pairs. Put
$s=\rho_i$ and let $b\in B$ be the other crossing endpoint.

\begin{lemma}
The stated configuration is impossible if $k\geq6$ and
$l\geq4$.
\end{lemma}
\begin{proof}
The two independent block-image generators must be the
literal neighbors of $i$, since every other generator fixes
both crossing endpoints. On $B$ they cannot both exchange
the blocks: they would then act identically on each of its
paired rows. Their distant neighbors, each commuting with
the opposite swap, and every other generator would fix the
point opposite $b$. Its pair would be $G_i$-invariant.
After reversal, therefore, write $\tau_L=\rho_{i-1}$,
$\tau_R=\rho_{i+1}$, where $\tau_L$ fixes $B$ pointwise
and $\tau_R$ exchanges its blocks. Independence forces
$\tau_L$ to exchange the blocks of $A$. The generator
$\tau_R$ may either fix or exchange those blocks; both
literal possibilities will be considered.

Both $k,l$ are even. To see this, choose in the group $M$
the element whose block image exchanges only $A$.
On $D$ it has two transpositions, and on the other paired
rows of $A$ it has $k-2$ transpositions; its sign is
$(-1)^k$. The element whose block image exchanges only
$B$ has two transpositions on $D$ and $l$ on $B$,
and sign $(-1)^l$. Both elements belong to $A_n$.
This parity argument does not replace either distinguished
generator in the marking.

Let $x=\rho_{i-2}$ and $y=\rho_{i+2}$ when present, and
let $C$ be all generators except these, $\tau_L,s,\tau_R$.
They centralize $M$. Since they fix $a$, they fix the whole
regular $M$-orbit $D$ pointwise. Their action on the other
paired rows is diagonal; they also fix the row containing
$b$. Thus their faithful row group is contained in
\[
 S_{k-2}\times S_{l-1}.
\]
Its independent length is at most $K-5$, where $K=k+l$.
The exact deficit formula gives $r=K$, so there are at
least $K-5$ generators in $C$, with equality only if both
$x,y$ exist. Consequently both exist and the equality
clause gives the full row product above, with a maximal
independent marking.

The prefix $P$ of $C$ and its suffix $Q$ commute and have
trivial intersection by IP. The two symmetric factors are
centerless and directly indecomposable, so they distribute
as whole factors between these intervals. The factor
$S_{l-1}$ must belong to $Q$: if it belonged to $P$,
$y$ would centralize its whole natural action, fixing the
opposite point to $b$; all other generators except $\tau_R$
fix it as well, and $\tau_R$ only exchanges the anchor
pair. This contradicts transitivity on $B$.

The other factor cannot also lie in $Q$. In that case $x$
centralizes the entire row group, and so fixes $B$ pointwise
and all points of $A\setminus D$. On $D$, it preserves
the two blocks, fixes $a$, and centralizes $\tau_R$.
If $\tau_R$ exchanges the blocks of $A$, this also fixes
its opposite point and hence all four points of $D$.
If $\tau_R$ fixes the blocks, its two transpositions on
$D$ make $x$ either the identity or the single
transposition on the block not containing $a$.
The first alternative makes $x$ trivial, the second odd.
Thus $P=S_{k-2}$ and $Q=S_{l-1}$.

The generator $y$ centralizes $P$ and $\tau_L$, preserves
both blocks of $A$, and fixes $a$. It fixes every point
outside $D$ because the corresponding natural symmetric
centralizers are trivial. It also fixes $a^{\tau_L}$,
and hence all four points of $D$. Therefore $y$ is the
identity on $A$.

If $\tau_R$ fixes the blocks of $A$, its support on $A$
is exactly $D$. The generator $x$ commutes with it and
therefore preserves $D$. The same is already true of
$M,C,y$, contradicting transitivity on $A$, since $k\geq6$.
It remains to consider the mixed basis, in which $\tau_R$
also exchanges the blocks of $A$.

If $l\geq6$, the maximal independent string of the natural
$Q=S_{l-1}$ is ordinary Coxeter; the outer degree-six
exception cannot occur because $l-1$ is odd. Its near-end
omission is $S_{l-2}$ fixing the anchor and one additional
row. Since $y$ centralizes this omission, fixes $b$, and
must connect the opposite anchor to another row, its action
on $B$ is one transposition. It is trivial on $A$ and so
odd, impossible. Thus only $l=4$ remains.

Write $p,q$ for the far and near generators of $Q$, so
they occur in the order $s,\tau_R,y,q,p$. In row labels
on $B$ they are
\[
 p=(1\ 2)(1'\ 2'),\qquad q=(2\ 3)(2'\ 3').
\]
Since $y$ is trivial on $A$, the accepted four-row
centralizer calculation gives exactly
$y=(1\ 2)(0'\ 3')$ or $(1'\ 2')(0'\ 3')$.
Moreover $s$ commutes with $Q,y$ and therefore fixes
$B\setminus\{b\}$ pointwise.

Consider the endpoint omission of the maximal prefix
marking of $P=S_{k-2}$. The accepted classifications,
including the elementary degree-four cases, make this
either a natural $S_{k-3}$ on all but one of its rows,
or a transitive subgroup on all $k-2$ rows. The latter
includes the twisted $S_4$ omission $D_8$ and the outer
$S_6$ omission. Since $x$ centralizes this omission,
that latter alternative would make $x$ preserve
$A\setminus D$, leaving $D$ invariant under $G_i$.
It is impossible. In the ordinary omission case, its
centralizer on either large block is the symmetric group
on the two rows of $D$ and its one omitted row.
The generator $x$ fixes $a$ and must connect $D$ to the
other rows, so on that block it is the transposition
between the other row of $D$ and the omitted row.
Its action on the opposite block is determined by its
commutation with $\tau_R$.

It follows that $\gen{P,x}$ is the natural diagonal
$S_{k-1}$ on two orbits matched by $\tau_R$, fixing
exactly $a,a^{\tau_R}$ in $A$. Since $s$ centralizes
this group and sends $a$ to $b$, it fixes $a^{\tau_R}$.
On the two remaining $(k-1)$-point orbits it either fixes
every point or exchanges corresponding points by
$\tau_R$. The first option would make $s$ the odd
transposition $(a\ b)$, so the second is forced.

This is precisely the literal action used in the accepted
four-row endpoint identity: the swap is $\tau_R$, its
four-row orbit is $B$, its even endpoint generator is $y$,
and $s$ exchanges all row pairs of $A$ except the anchor
pair, replacing that exchange by $(a\ b)$. Therefore
\[
                   p=s(\tau_R s y)^5\tau_R y.
\]
The accepted ten-point cycle proof and cancelling extra
row pairs apply without modification. It makes the
distinguished generator $p$ redundant, the final contradiction.
\end{proof}

\begin{lemma}[A small central factor]
Let $m\geq3$ be odd. An independent involutory string
C-marking of $C_2\times S_m$ with $m$ generators contains
the central involution as a distinguished generator.
The remaining generators form the graph of a character
$S_m\to C_2$, with a maximal symmetric string marking.
\end{lemma}
\begin{proof}
For $m\geq5$, choose a minimal subfamily generating the
$S_m$ quotient. The accepted normal-quotient independent-length
bound shows that exactly $m-1$ generators are needed, since
the kernel contributes at most one. Their quotient images
are an extremal involutory string, hence an ordinary
Coxeter path by the accepted L047 classification; the
even-degree outer exception is absent. Its noncommuting
neighbors force the chosen subfamily to occupy consecutive
indices in the original string. The sole remaining
generator centralizes all but possibly one endpoint of
this path. The endpoint omission's centralizer in $S_m$
is trivial, so that generator is the central $C_2$
involution. IP makes the other subgroup intersect it
trivially, proving the character-graph statement. This is
the accepted small-factor proof, with its degree-five
endpoint supplied by the same accepted classification.

For $m=3$, write the central involution as $z$. All other
involutions are $z^\epsilon t$ with $t$ a transposition.
If no distinguished generator is $z$, the commuting
string ends have the same transposition image. Being
distinct, they are $t$ and $zt$. The middle generator
has a different transposition image. One of its two
products with the ends has order six, so that pair
already generates $C_2\times S_3$, a contradiction.
Thus $z$ occurs. The other two generators project to
distinct transpositions; their product must have order
three, since otherwise they already generate the whole
group. They therefore form the asserted complement.
\end{proof}

\begin{lemma}[The remaining anchored block sizes]
The regular-$M$-orbit configuration in this note is
impossible for all $k,l\geq2$ with $n\geq12$.
\end{lemma}
\begin{proof}
The parity argument above still gives even $k,l$.
The paired anchor orbit in $B$ still forces $y$ to
exist and to move the point opposite $b$.
If $l=2$, its action on $B$ is consequently a single
transposition: it fixes the first two-point block and
exchanges the second. On $A$, it preserves the blocks
and commutes with $\tau_L$, which exchanges them.
Its action there is diagonal and even. Thus $y$ is odd,
impossible. Hence take $l\geq4$.

Suppose $k=2$, so $A=D$ has four points. The generator
$x$, if present, is the identity on $A$. Indeed, if
$\tau_R$ exchanges its blocks, commutation and fixation
of $a$ fixes one point in each two-point block and thus
all four points. If $\tau_R$ fixes those blocks, $x$
can at most exchange the second block; on $B$ it is
diagonal by commutation with $\tau_R$, so evenness again
forces its four-point action to be the identity.
The whole left interval, consisting of $x$ and all
earlier generators, is therefore supported on $B$ and
fixes $b$. It commutes with the other generators of
$G_i$: the only possible exception in the string is
$\tau_L$, whose support is $A$, so they commute in fact.
Since $G_i$ is transitive on $B$, this left interval
is trivial. In detail its commuting complementary
factor has the whole orbit of $b$, so any element of
the left factor fixes every point in that orbit.
Consequently $x$ and the left interval are absent.

Counting now gives a maximal right interval $Q=S_{l-1}$
on the nonanchor rows of $B$. The generator $y$ is
trivial on $A$, since it centralizes $\tau_L$, preserves
the two blocks and fixes $a$. For $l\geq6$ the ordinary
endpoint omission again makes it odd. For $l=4$, its
two even possibilities are exactly the four-row
endpoint alternatives, and $s$ fixes all seven
nonanchor points of $B$. On $A\cup\{b\}$ it is the
crossing transposition $(a\ b)$ times one transposition
on the three remaining points of $A$. Choose coordinates
so that the nonidentity involution $\tau_R|_A$ is
$(A_0\ A'_0)(A_1\ A'_1)$ and $a=A_0$.
These are precisely the six cases already proved in
the accepted two-row endpoint: its words $w=p$ or
$(\tau_R s\tau_R w)^3=p$ make the far generator of
$Q$ redundant. No change of distinguished generators
is involved in this coordinate relabeling.

Finally suppose $k=4$. The independent-length count
used above again forces both $x,y$ to exist and gives
the full row group
\[
             \gen C=C_2\times S_{l-1}.
\]
Its central involution $z$ exchanges the two unused
rows of $A$ diagonally; its support is exactly
$A\setminus D$. Put $m=l-1$, which is odd. The
preceding lemma applies to the maximal rank-$m$
marking of $C$.

The noncentral $S_m$ component must lie in the right
interval: otherwise $y$ centralizes its entire natural
action on each block of $B$, fixing the opposite
anchor. If the central generator $z$ also lies in the
right interval, its final generator adjacent to $y$
must be in the $S_m$ component, for the same reason.
Thus $y$ commutes with $z$, while $x$ centralizes all
of $C$. Both preserve its support $A\setminus D$;
the remaining generators preserve $D$ as well, a
contradiction. Therefore the left interval is the
single distinguished generator $z$, and the right
interval is a character graph over $S_m$.

The character graph is trivial. If it were nontrivial,
commutation of $x$ with an odd graph element would
force $x|_A$ to centralize $z|_A$ (restrict the
commutator to $A$). Hence $x$ would preserve $D$.
Also $y$ commutes with $z$ and so preserves $D$,
contradicting transitivity on $A$ again.
We thus have a pure right-hand $S_m$ factor.

The generator $y$ centralizes $z$ and $\tau_L$,
preserves the blocks and fixes $a$. It fixes $D$
pointwise, and on its complement is either the
identity or $z$. In particular its action on $A$
is even. For $l\geq6$ its action on $B$ is the
single odd endpoint transposition, again impossible.
Take $l=4$, and write $y|_A=z^\eta$ with
$\eta\in\{0,1\}$. The element $y_0=yz^\eta$ is
supported on $B$, with one of the two even four-row
endpoint actions. It is a word in available generators,
not a replacement marking. The generator $s$ commutes
with $z,y$ and the pure right interval, hence with
$y_0$; it fixes $B\setminus\{b\}$ pointwise.

The pure-block possibility for $\tau_R|_A$ is
impossible, since its support $D$ would be preserved
by $x$, which centralizes it. In the mixed possibility
$\tau_R$ exchanges the blocks of $A$. The generator
$x$ fixes $a$ and centralizes this swap. On the
three other points of its first block it must exchange
the remaining point of $D$ with one of the unused
rows, to connect $D$ to its complement; this determines
its action on the opposite block. Hence $\gen{z,x}$
is a diagonal $S_3$ on two orbits matched by
$\tau_R$, fixing $a,a^{\tau_R}$.
Centralization and evenness now force $s$ to exchange
those two three-point orbits by $\tau_R$, fixing
$a^{\tau_R}$ and exchanging $a,b$. The accepted
four-row identity therefore gives
\[
                   p=s(\tau_R s y_0)^5\tau_R y_0,
\]
where $p$ is the far generator of the right interval.
Its right side uses only distinguished generators
other than $p$, making $p$ redundant. This proves
the final small case.
\end{proof}

\paragraph{Unfinished scope.}
This proof addresses the case where the regular
$M$-orbit contains its crossing endpoint. Regular
orbits away from that endpoint remain separate, as
does the $\beta=0$, $\gamma=1$ branch and the other
source obligations. The pure and mixed marked bases
are distinguished throughout. This is not a whole
L044 recovery or a new-question claim.
\paragraph{Attribution and paid dependencies.}
The source problem and orbit-cut framework are from Cameron, Fernandes,
Leemans and Mixer, \emph{Highest rank of a polytope for $A_n$},
arXiv:1605.09173v1, Section~5.1. The present marked-basis and endpoint
arguments are new proofs of bounded cases in that framework. They use
the complete accepted block-deficit, small-endpoint and two-row notes,
and the full retained L047 extremal/normal-quotient interfaces, including
Whiston's independent-generation bounds. Their exact files and accepted
reviews are bound in the companion dependency map; no sibling proof is
an uncharged premise. No new finite census or formal theorem is used.
\end{document}
